\documentclass[
    language=chinese,
    doctype=article,
    institution=none,
]{../../omnilatex}

\RequirePackage{config/document-settings}

\begin{document}

\maketitle

\section*{摘要}
本文档展示了 OmniLaTeX 对中文排版的全面支持，包括字体配置、行距处理和拼音标注。同时，本文以软件工程为主题，讨论中文语境下的技术写作实践。

\tableofcontents

\section{引言}
OmniLaTeX 基于 LuaLaTeX 和 luatexja 引擎，为中文文档提供原生排版支持。
系统使用 Noto Serif CJK SC 作为默认中文字体，确保在各种操作系统上的一致显示效果。

在当今全球化的软件开发环境中，中文技术文档的质量直接影响到团队协作效率和产品本地化的成功率。一份排版规范、格式统一的技术文档，不仅降低了读者的认知负担，也为跨国团队的沟通搭建了桥梁。

\section{中文排版特性}

\subsection{基本排版}
天地玄黄，宇宙洪荒。日月盈昃，辰宿列张。寒来暑往，秋收冬藏。
闰余成岁，律吕调阳。云腾致雨，露结为霜。金生丽水，玉出昆冈。
剑号巨阙，珠称夜光。果珍李柰，菜重芥姜。

传统中文排版遵循严格的版式规则：段落首行缩进两个字符宽度，行距一般为字号的 1.5 至 2 倍，标点符号使用全角形式。OmniLaTeX 自动处理这些细节，无需手动调整。

\subsection{现代汉语与中英混排}
中文与英文混排时，OmniLaTeX 会自动在中文和 Latin 字符之间插入适当的间距。
例如：使用 LuaLaTeX 编译引擎和 luatexja 宏包可以实现高质量的中文排版效果。

在技术文档中，中英混排是常态。算法描述中常见的术语如「时间复杂度为 $O(n \log n)$」「使用 HTTP/2 协议进行通信」等场景，都需要在中英文之间保持自然、美观的间距。OmniLaTeX 的间距处理机制确保这些混合内容不会出现拥挤或分离过度的问题。

\subsection{拼音标注}
OmniLaTeX 支持通过 ruby 标注为汉字添加拼音：

%pinyin{汉}{hàn}\pinyin{字}{zì}\pinyin{拼}{pīn}\pinyin{音}{yīn}

拼音标注在汉语教学和古籍数字化领域有广泛应用。OmniLaTeX 的 ruby 注解机制与 Unicode 标准完全兼容，支持声调符号的正确渲染。

\subsection{列表与枚举}
中文技术文档经常使用有序和无序列表来组织内容。OmniLaTeX 对列表项的缩进、标记符号和行距进行了本地化适配，确保列表在中文环境下的视觉效果与中文排版规范一致。

\section{软件工程中的中文技术写作}

\subsection{接口文档}
软件接口文档是开发者最常阅读的技术文档类型。一份优秀的中文接口文档应当包含以下要素：接口的 URL 路径、请求方法、请求参数及类型、返回值结构、错误码说明，以及至少一个调用示例。

例如，一个用户认证接口的文档可以这样组织：首先描述接口的用途和适用场景，然后列出请求参数表格，最后给出 JSON 格式的请求和响应示例。使用 OmniLaTeX 的表格环境和代码高亮宏包，可以轻松实现专业级别的文档排版。

\subsection{算法文档}
算法文档是中文技术写作中的另一个重要类别。与英文算法文档类似，中文算法文档也需要使用数学公式来精确描述算法的时间复杂度和空间复杂度。例如，快速排序的平均时间复杂度为 $O(n \log n)$，最坏情况下为 $O(n^2)$。

OmniLaTeX 完整支持 \LaTeX{} 的数学排版功能，包括行内公式、行间公式、矩阵、多行对齐等。这使得中文算法文档可以使用与英文文献相同的数学表达规范。

\subsection{开源项目文档}
随着中国开源社区的发展，越来越多的开源项目选择提供中文文档。OmniLaTeX 可用于排版项目说明、贡献指南、架构设计文档和发布说明等多种文档类型。统一的排版风格有助于提升项目的专业形象和社区参与度。

\subsection{数学公式示例}
中文数学文档经常需要排版复杂的公式。OmniLaTeX 完整支持 \LaTeX{} 的数学环境。例如，欧拉公式可以写作 $e^{i\pi} + 1 = 0$，而高斯积分可以表示为：
\[
\int_{-\infty}^{+\infty} e^{-x^2} \, dx = \sqrt{\pi}
\]

\section{结论}
OmniLaTeX 为中文学术文档提供了完整的排版解决方案。通过 LuaLaTeX 和 luatexja 的深度集成，它实现了对简体中文、繁体中文以及中英混排内容的原生支持。无论是学术论文、技术报告还是软件文档，OmniLaTeX 都能提供一致、高质量的排版效果。

\end{document}
